\section*{Примечания о чтениях}
\subsection*{\texttt{GELFAND-P001405}}
\par\noindent\leavevmode\phantomsection\label{GELFAND-0153-U001}\textbf{\texttt{GELFAND-0153-U001}} (pairing-value).\nopagebreak[4]\par
В формуле (35) после знака равенства виден открытый справа С-образный знак. Bounding box original pixels: [720,125,890,280].\par
Кандидаты: 0; C. Чтение не установлено в дипломатическом исходнике.\par
Историческое примечание исходной транскрипции: Область просмотрена при увеличении 10x. Знак может быть незамкнутым нулём либо C; математически ожидаемое значение не подставлялось. Соседние индексы и условие J ≠ I′ сохранены.\par\medskip
\par\noindent\leavevmode\phantomsection\label{GELFAND-0153-U002}\textbf{\texttt{GELFAND-0153-U002}} (characters-line-ending).\nopagebreak[4]\par
В конце строки после «характеры» видны с-, затем несколько слитных скорописных знаков и конечный дефис; следующая строка начинается «-ф-и Шуберта». Bounding box original pixels: [1980,1020,2360,1120].\par
Кандидаты: не записаны. Чтение не установлено в дипломатическом исходнике.\par
Историческое примечание исходной транскрипции: Просмотрено при 3x и 6x. Точный текст промежуточного сокращения не установлен; начальное с- и конечный дефис отражены в маркере, продолжение следующей строки сохранено.\par\medskip
\subsection*{\texttt{GELFAND-P001406}}
\par\noindent\leavevmode\phantomsection\label{GELFAND-0153-U101}\textbf{\texttt{GELFAND-0153-U101}} (interchange-diagram).\nopagebreak[4]\par
Диаграмма (1): два вертикально разделённых прямоугольника с A/B и C/густо обведённым D-подобным знаком, кружок композиции, равенство и справа блок с рядами A,C и B,D-подобный знак. В узкой средней полосе кружок; у боковых концов повторно обведённые стрелочные штрихи. Bounding box original pixels: [145,180,1150,440].\par
Кандидаты: не записаны. Чтение не установлено в дипломатическом исходнике.\par
Историческое примечание исходной транскрипции: Оригинал и 3x-кроп просмотрены. Направления отдельных штрихов справа и состояние переписанных букв не восстановлены догадкой; описание удерживает расположение уверенно видимых элементов.\par\medskip
\par\noindent\leavevmode\phantomsection\label{GELFAND-0153-U102}\textbf{\texttt{GELFAND-0153-U102}} (example-cancel-left).\nopagebreak[4]\par
После «не вып.» перед «акс.» зачёркнута небольшая группа знаков. Bounding box original pixels: [970,465,1050,530].\par
Кандидаты: не записаны. Чтение не установлено в дипломатическом исходнике.\par
Историческое примечание исходной транскрипции: Зачёркнутый текст не читается надёжно после просмотра 4x-кропа; факт отмены сохранён.\par\medskip
\par\noindent\leavevmode\phantomsection\label{GELFAND-0153-U103}\textbf{\texttt{GELFAND-0153-U103}} (example-cancel-right).\nopagebreak[4]\par
После «акс.» перед незачёркнутым «(1)» густо зачёркнут знак или прежняя ссылка в скобках. Bounding box original pixels: [1165,460,1300,565].\par
Кандидаты: не записаны. Чтение не установлено в дипломатическом исходнике.\par
Историческое примечание исходной транскрипции: Область просмотрена при 4x. Состав зачёркнутой группы нельзя установить уверенно.\par\medskip
\par\noindent\leavevmode\phantomsection\label{GELFAND-0153-U104}\textbf{\texttt{GELFAND-0153-U104}} (tensor-indices).\nopagebreak[4]\par
Формула по расположению знаков читается как (AB)\_\{ik\}\textasciicircum{}\{rs\}=Σ A\_\{ij\}\textasciicircum{}\{kℓ\} B\_\{js\}\textasciicircum{}\{ℓr\}. В углах A стоят k,ℓ сверху и i,j снизу; у B — ℓ,r сверху и j,s снизу. Bounding boxes original pixels: formula [390,950,1230,1160], tensor boxes [380,690,950,910].\par
Кандидаты: (AB)\_\{ik\}\textasciicircum{}\{rs\}=Σ A\_\{ij\}\textasciicircum{}\{kℓ\} B\_\{js\}\textasciicircum{}\{ℓr\}. Чтение не установлено в дипломатическом исходнике.\par
Историческое примечание исходной транскрипции: Индексы формулы просмотрены при 6x и 10x. Положение свободных индексов слева не согласуется с их положением справа; возможная ошибка источника сохранена, не исправлена. Буква r начертана τ-образно.\par\medskip
\par\noindent\leavevmode\phantomsection\label{GELFAND-0153-U105}\textbf{\texttt{GELFAND-0153-U105}} (boundary-diagram).\nopagebreak[4]\par
В небольшом квадрате видны Δu=0, горизонтальные стрелки вправо и вертикальные вверх; слева u и ∂u/∂n. Над левой подписью двухчастный знак с открытым C-подобным началом и высоким вторым штрихом. Bounding box original pixels: [245,2260,550,2455]; ambiguous symbol [248,2290,360,2380].\par
Кандидаты: не записаны. Чтение не установлено в дипломатическом исходнике.\par
Историческое примечание исходной транскрипции: Квадрат и подписи просмотрены при 6x, неясный знак — при 10x. Его точное чтение и сопряжение внутренних стрелок с границей остаются неустановленными; диаграмма сохранена через ссылку на источник.\par\medskip
\subsection*{\texttt{GELFAND-P001407}}
\par\noindent\leavevmode\phantomsection\label{GELFAND-0153-U201}\textbf{\texttt{GELFAND-0153-U201}} (last-line-cancellation).\nopagebreak[4]\par
В последней строке после «были определены» и перед «ещё» зачёркнуто короткое слово или группа знаков. Bounding box original pixels: [1970,2925,2085,2990].\par
Кандидаты: не записаны. Чтение не установлено в дипломатическом исходнике.\par
Историческое примечание исходной транскрипции: Просмотрено при 6x и 10x. Под несколькими перекрывающими штрихами не удаётся надёжно прочесть отменённый текст; незачёркнутое «ещё» сохранено без редакционного исправления.\par\medskip
\subsection*{\texttt{GELFAND-P001408}}
\par\noindent\leavevmode\phantomsection\label{GELFAND-0153-U301}\textbf{\texttt{GELFAND-0153-U301}} (cancelled-title).\nopagebreak[4]\par
Между именем «И. Захаревич» и заголовком имеется слово или короткая запись, густо зачёркнутая горизонтальными и наклонными линиями.\par
Кандидаты: неразборчиво. Чтение не установлено в дипломатическом исходнике.\par
Историческое примечание исходной транскрипции: Исходная область bbox [595,85,840,178], координаты x0,y0,x1,y1, пиксели. Проверена в оригинале и с увеличением 6x. Начальные и конечные штрихи не дают надёжного полного чтения.\par\medskip
\par\noindent\leavevmode\phantomsection\label{GELFAND-0153-U302}\textbf{\texttt{GELFAND-0153-U302}} (title-preposition).\nopagebreak[4]\par
В заголовке «Импровизация … хаосе.» предлог имеет открытый округлый контур, похожий на о либо с.\par
Кандидаты: о; с. Чтение не установлено в дипломатическом исходнике.\par
Историческое примечание исходной транскрипции: Исходная область bbox [1255,75,1365,170]; контекст заголовка bbox [855,70,1570,185]. Проверено 6x и независимым повторным чтением. Грамматически ожидаемое «о» не подставлено молча.\par\medskip
\par\noindent\leavevmode\phantomsection\label{GELFAND-0153-U303}\textbf{\texttt{GELFAND-0153-U303}} (baker-cancellations).\nopagebreak[4]\par
После «этим» зачёркнута короткая запись; ниже имеются отдельный короткий зачёркнутый фрагмент и длинная зачёркнутая строка с буквоподобными и Z-подобными знаками. Над строкой разборчиво вставлено «(с точн. до меры 0)».\par
Кандидаты: не записаны. Чтение не установлено в дипломатическом исходнике.\par
Историческое примечание исходной транскрипции: Исходная область bbox [1490,850,2250,990]. Вся область и её увеличенная копия 6x просмотрены; точные слова и математические знаки под плотным зачёркиванием не восстановлены.\par\medskip
\par\noindent\leavevmode\phantomsection\label{GELFAND-0153-U304}\textbf{\texttt{GELFAND-0153-U304}} (baker-sequence).\nopagebreak[4]\par
Три квадратных рамки с изменяющимися округлыми и вытянутыми линиями; у первого рисунка есть горизонтальный пунктир, у третьего — отчётливая направленная влево нижняя горизонтальная линия. На иных мелких окончаниях нельзя уверенно отделить стрелки от утолщений пера.\par
Кандидаты: не записаны. Чтение не установлено в дипломатическом исходнике.\par
Историческое примечание исходной транскрипции: Исходная область bbox [425,1320,1480,1575]. Проверена в оригинале и в увеличенном 3x фрагменте. Сохранено описание всех трёх рисунков; точная геометрия и неуверенные направления не реконструированы.\par\medskip
\par\noindent\leavevmode\phantomsection\label{GELFAND-0153-U305}\textbf{\texttt{GELFAND-0153-U305}} (hyperbolic-surface).\nopagebreak[4]\par
Контур с двумя расширениями, двумя внутренними дугообразными фигурами и центральным наклонным отрезком; от середины отрезка идёт стрелка вправо-вверх. Направление некоторых коротких штрихов на дугах и точное соединение кривых неясны.\par
Кандидаты: не записаны. Чтение не установлено в дипломатическом исходнике.\par
Историческое примечание исходной транскрипции: Исходная область bbox [480,2440,1030,2690]. Осмотрено в оригинале и в увеличении 3x. Подпись о геодезических потоках вынесена в транскрипцию; рисунок оставлен для сверки с источником.\par\medskip
\par\noindent\leavevmode\phantomsection\label{GELFAND-0153-U306}\textbf{\texttt{GELFAND-0153-U306}} (manifold-abbreviation).\nopagebreak[4]\par
В нижней подписи после «геодезич. потоки на» сокращённое слово начинается «мн-», далее видны слитные выносные и округлые штрихи и окончание «иях».\par
Кандидаты: мн-зиях; мн-иях. Чтение не установлено в дипломатическом исходнике.\par
Историческое примечание исходной транскрипции: Исходная область bbox [1690,2400,1940,2490]; широкий контекст bbox [1090,2380,2400,2490]. Проверено в оригинале и при увеличении 6x. Сокращение не развёрнуто до ожидаемого полного слова.\par\medskip
\subsection*{\texttt{GELFAND-P001409}}
\par\noindent\leavevmode\phantomsection\label{GELFAND-0153-U401}\textbf{\texttt{GELFAND-0153-U401}} (degree-one-cancellation).\nopagebreak[4]\par
Между «deg P=1» и «неинтересный случай» присутствует короткое зачёркнутое слово или знак; над ним отдельная горизонтальная черта, рядом виден угловой штрих.\par
Кандидаты: не записаны. Чтение не установлено в дипломатическом исходнике.\par
Историческое примечание исходной транскрипции: Исходная область bbox [450,440,640,555]. Проверена при увеличении 6x. Под густой штриховкой полное чтение не установлено.\par\medskip
\par\noindent\leavevmode\phantomsection\label{GELFAND-0153-U402}\textbf{\texttt{GELFAND-0153-U402}} (unit-modulus-claim).\nopagebreak[4]\par
В источнике непосредственно написано «|a|=1 ⇒ всё разбив. на инвар. окружности, на к-рых поворот.»; дополнительного ограничения на a в этой строке не видно.\par
Кандидаты: |a|=1 ⇒ всё разбив. на инвар. окружности, на к-рых поворот.. Чтение не установлено в дипломатическом исходнике.\par
Историческое примечание исходной транскрипции: Исходная область bbox [230,900,2165,985]. Формулировка разборчива и сохранена. Математическая оговорка: при a=1 и b≠0 отображение z↦az+b является переносом; исключение a=1 здесь не записано. Исправление в исходный текст не внесено.\par\medskip
\subsection*{\texttt{GELFAND-P001410}}
\par\noindent\leavevmode\phantomsection\label{GELFAND-0153-U501}\textbf{\texttt{GELFAND-0153-U501}} (isolated-cancellation).\nopagebreak[4]\par
Под первой строкой определения, слева, находится короткий пучок перечёркнутых наклонных штрихов; первоначальная запись могла содержать U, но не читается уверенно.\par
Кандидаты: U. Чтение не установлено в дипломатическом исходнике.\par
Историческое примечание исходной транскрипции: Исходная область bbox [255,390,405,495]. Проверена в оригинале и при увеличении 6x. Эта самостоятельная зачёркнутая запись учтена отдельно.\par\medskip
\par\noindent\leavevmode\phantomsection\label{GELFAND-0153-U502}\textbf{\texttt{GELFAND-0153-U502}} (normality-cancellation).\nopagebreak[4]\par
В строке о точках x после «∃ окр. U —» и перед «окр. нормальности» короткое слово густо зачёркнуто.\par
Кандидаты: не записаны. Чтение не установлено в дипломатическом исходнике.\par
Историческое примечание исходной транскрипции: Исходная область bbox [1800,535,1940,625]. Просмотрено при увеличении 6x; сохранившиеся пересекающиеся штрихи не позволяют уверенного чтения.\par\medskip
\par\noindent\leavevmode\phantomsection\label{GELFAND-0153-U503}\textbf{\texttt{GELFAND-0153-U503}} (components-cancellation).\nopagebreak[4]\par
Под словами «мн-ва нормальности» имеется зачёркнутая строка с несколькими начальными и конечными петлями; центральные буквы закрыты несколькими горизонтальными линиями.\par
Кандидаты: не записаны. Чтение не установлено в дипломатическом исходнике.\par
Историческое примечание исходной транскрипции: Исходная область bbox [1490,1170,1900,1250]. Проверена 6x. Полное содержание зачёркнутой строки не восстановлено.\par\medskip
\par\noindent\leavevmode\phantomsection\label{GELFAND-0153-U504}\textbf{\texttt{GELFAND-0153-U504}} (sullivan-cancellation).\nopagebreak[4]\par
Перед «преобр. дискр. мн-ва» виден зачёркнутый фрагмент из петель и наклонных штрихов; крайний левый знак похож на ф либо f.\par
Кандидаты: ф; f. Чтение не установлено в дипломатическом исходнике.\par
Историческое примечание исходной транскрипции: Исходная область bbox [285,1400,485,1520]. Проверена 6x. Ни весь фрагмент, ни его точное число знаков не установлены.\par\medskip
\par\noindent\leavevmode\phantomsection\label{GELFAND-0153-U505}\textbf{\texttt{GELFAND-0153-U505}} (structure-cancellation).\nopagebreak[4]\par
Между «конф.» и «структуру» короткое зачёркнутое слово; видны нижний вынос в начале и верхняя петля ближе к концу.\par
Кандидаты: преобр.. Чтение не установлено в дипломатическом исходнике.\par
Историческое примечание исходной транскрипции: Исходная область bbox [955,2120,1180,2225]. Просмотрено 6x. Кандидат основан на видимых остатках букв и не утверждается. Следующее слово «структуру» сохранено именно в видимой форме.\par\medskip
\par\noindent\leavevmode\phantomsection\label{GELFAND-0153-U506}\textbf{\texttt{GELFAND-0153-U506}} (circle-packing).\nopagebreak[4]\par
Квадрат с четырьмя крупными почти круглыми контурами и малыми замкнутыми кривыми в промежутках, у сторон и углов; справа стрелка влево с подписью «вписываем». Небольшие фигуры и касания местами сливаются с жирными штрихами рамки.\par
Кандидаты: не записаны. Чтение не установлено в дипломатическом исходнике.\par
Историческое примечание исходной транскрипции: Исходная область рисунка и подписи bbox [300,2370,1150,2695]; плотная геометрия bbox [300,2370,670,2695] осмотрена с увеличением 4x. Не подставлена идеальная или математически достроенная упаковка.\par\medskip
\subsection*{\texttt{GELFAND-P001411}}
\par\noindent\leavevmode\phantomsection\label{GELFAND-0153-U601}\textbf{\texttt{GELFAND-0153-U601}} (example-cancelled).\nopagebreak[4]\par
В исходном прямоугольнике (470,605,620,700) над рисунком справа от «Пример» расположен короткий фрагмент, перечёркнутый несколькими наклонными линиями.\par
Кандидаты: не записаны. Чтение не установлено в дипломатическом исходнике.\par
Историческое примечание исходной транскрипции: Проверено по оригиналу и при увеличении 8×. Под зачёркиванием остаются дугообразные и вертикальные штрихи; буквенное или математическое чтение не установлено.\par\medskip
\par\noindent\leavevmode\phantomsection\label{GELFAND-0153-U602}\textbf{\texttt{GELFAND-0153-U602}} (book-after-rota).\nopagebreak[4]\par
В исходном прямоугольнике (815,1830,975,1945), после «Рота», виден изогнутый знак, похожий на э или закрывающую скобку, и две отдельные точки справа. Нижний край знака соприкасается с подчёркиванием.\par
Кандидаты: э..; )... Чтение не установлено в дипломатическом исходнике.\par
Историческое примечание исходной транскрипции: Проверено при увеличении 7×. Неясно, является ли изогнутый штрих текстовым знаком или завитком линии подчёркивания. В body.tex оставлено возможное чтение без нормализации.\par\medskip
\par\noindent\leavevmode\phantomsection\label{GELFAND-0153-U603}\textbf{\texttt{GELFAND-0153-U603}} (book-overwritten-noun).\nopagebreak[4]\par
В исходном прямоугольнике (1130,1810,1330,1960) в приписке «см. … Айгнера (?)» начало слова густо перезаписано; правые буквы выглядят как «ниге».\par
Кандидаты: книге; книгу; книжке. Чтение не установлено в дипломатическом исходнике.\par
Историческое примечание исходной транскрипции: Проверено при увеличении 5× и независимо вторым читателем. Видимые штрихи не позволяют надёжно прочитать весь фрагмент или восстановить зачёркнутый первоначальный вариант. Написанное «(?)» после фамилии воспроизведено как авторский знак вопроса.\par\medskip
\par\noindent\leavevmode\phantomsection\label{GELFAND-0153-U604}\textbf{\texttt{GELFAND-0153-U604}} (basis-exchange-cancelled).\nopagebreak[4]\par
В исходном прямоугольнике (1470,2315,1650,2450), после (B\_1\textbackslash{}x)∪y и перед запятой, короткий фрагмент зачёркнут плотными наклонными штрихами.\par
Кандидаты: не записаны. Чтение не установлено в дипломатическом исходнике.\par
Историческое примечание исходной транскрипции: Проверено при увеличении 8×. Различимы остатки связных петель, но чтение букв или символов ненадёжно.\par\medskip
\par\noindent\leavevmode\phantomsection\label{GELFAND-0153-U605}\textbf{\texttt{GELFAND-0153-U605}} (problem-four-cancelled).\nopagebreak[4]\par
В исходном прямоугольнике (220,3230,330,3345), между номером (4) и «Задача», расположен узкий зачёркнутый фрагмент.\par
Кандидаты: не записаны. Чтение не установлено в дипломатическом исходнике.\par
Историческое примечание исходной транскрипции: Проверено при увеличении 8×. Исходные знаки не отделяются от штрихов зачёркивания.\par\medskip
\subsection*{\texttt{GELFAND-P001412}}
\par\noindent\leavevmode\phantomsection\label{GELFAND-0153-U701}\textbf{\texttt{GELFAND-0153-U701}} (top-cancelled).\nopagebreak[4]\par
В исходном прямоугольнике (1320,60,1415,160), между «из» и «Гельфанда», виден короткий фрагмент с густо перезаписанными вертикальными и петлеобразными штрихами.\par
Кандидаты: не записаны. Чтение не установлено в дипломатическом исходнике.\par
Историческое примечание исходной транскрипции: Проверено по оригиналу и при увеличении 8×; надёжное чтение первоначальных знаков невозможно.\par\medskip
\par\noindent\leavevmode\phantomsection\label{GELFAND-0153-U702}\textbf{\texttt{GELFAND-0153-U702}} (knuth-name).\nopagebreak[4]\par
В исходном прямоугольнике (550,160,875,285), после «связано с», фамилия по видимому порядку штрихов напоминает «Кунтом». Средние буквы слиты.\par
Кандидаты: Кунтом; Кнутом. Чтение не установлено в дипломатическом исходнике.\par
Историческое примечание исходной транскрипции: Проверено при увеличении 8×. Возможная перестановка букв в источнике не исправлена по ожидаемой фамилии.\par\medskip
\par\noindent\leavevmode\phantomsection\label{GELFAND-0153-U703}\textbf{\texttt{GELFAND-0153-U703}} (determinant-formula-source).\nopagebreak[4]\par
В исходном прямоугольнике (275,270,2320,425) формула (6) содержит сумму без видимых пределов, множитель (-1)\textasciicircum{}\{d-i\}, первый правый детерминант [y\_1,x\_2,…,x\_d] и второй правый детерминант [y\_1,…,y\_\{i-1\},x\_i,y\_\{i+1\},…,y\_d].\par
Кандидаты: y\_1 в первом правом детерминанте; x\_i во втором; (-1)\textasciicircum{}\{d-i\} перед произведением. Чтение не установлено в дипломатическом исходнике.\par
Историческое примечание исходной транскрипции: Индексы проверены при увеличении 6× и 8×. Первая нижняя цифра сходна с другими единицами на странице; у x во втором детерминанте видна точка над i. Сочетание индексов и знака математически требует проверки; рукопись не заменена предполагаемой стандартной формулой.\par\medskip
\par\noindent\leavevmode\phantomsection\label{GELFAND-0153-U704}\textbf{\texttt{GELFAND-0153-U704}} (overwritten-eight).\nopagebreak[4]\par
В исходном прямоугольнике (130,900,265,1020) номер пункта в круглых скобках густо перезаписан; наружный контур похож на 8.\par
Кандидаты: 8. Чтение не установлено в дипломатическом исходнике.\par
Историческое примечание исходной транскрипции: Проверено при увеличении 8×. Окончательное чтение выглядит как 8, но первоначальная цифра и порядок наложенных штрихов не установлены.\par\medskip
\par\noindent\leavevmode\phantomsection\label{GELFAND-0153-U705}\textbf{\texttt{GELFAND-0153-U705}} (problem-eleven-cancelled).\nopagebreak[4]\par
В исходном прямоугольнике (1800,1595,1940,1690), после стрелки и перед «2 матроида», короткое слово или знак перечёркнут горизонтальными линиями.\par
Кандидаты: не записаны. Чтение не установлено в дипломатическом исходнике.\par
Историческое примечание исходной транскрипции: Проверено при увеличении 8×. Видимые остатки не позволяют уверенно восстановить зачёркнутое чтение.\par\medskip
\par\noindent\leavevmode\phantomsection\label{GELFAND-0153-U706}\textbf{\texttt{GELFAND-0153-U706}} (basis-construction-word).\nopagebreak[4]\par
В исходном прямоугольнике (720,1890,1190,2040) справа от стрелки к решётке написано «по опр.» и подчёркнутое слово «баз…». Последняя буква неоднозначна.\par
Кандидаты: база; базы; базе. Чтение не установлено в дипломатическом исходнике.\par
Историческое примечание исходной транскрипции: Проверено при увеличении 7× и независимо вторым читателем. Наиболее вероятно буквальное «база», но последняя буква остаётся неоднозначной. Слово не изменено по грамматическому или математическому ожиданию.\par\medskip
\par\noindent\leavevmode\phantomsection\label{GELFAND-0153-U707}\textbf{\texttt{GELFAND-0153-U707}} (basis-grid).\nopagebreak[4]\par
В исходном прямоугольнике (200,1890,760,2140) изображены четыре вертикальные и четыре горизонтальные пересекающиеся линии. Слева и снизу фигурные скобки; слева подпись похожа на «баз», снизу — на «база» или «базы». Справа горизонтальная стрелка направлена влево к решётке.\par
Кандидаты: левая подпись: баз; нижняя подпись: база; левая подпись: баз; нижняя подпись: базы. Чтение не установлено в дипломатическом исходнике.\par
Историческое примечание исходной транскрипции: Проверено при увеличении 8×. Геометрия различима, но точное чтение подписей не установлено; оставлена ссылка на исходный рисунок без выдуманного уточнения подписей.\par\medskip
\par\noindent\leavevmode\phantomsection\label{GELFAND-0153-U708}\textbf{\texttt{GELFAND-0153-U708}} (basis-axiom-reference).\nopagebreak[4]\par
В исходном прямоугольнике (2100,695,2200,810), в конце замечания (7), после «аксиому замены» отчётливо написано «(3)».\par
Кандидаты: (3). Чтение не установлено в дипломатическом исходнике.\par
Историческое примечание исходной транскрипции: Ссылка сохранена по источнику. На предыдущем кадре аксиоматика отмечена номером (2), а (3) начинает обсуждение биматроида; возможное несоответствие внутренней ссылки не исправлено.\par\medskip
\subsection*{\texttt{GELFAND-P001413}}
\par\noindent\leavevmode\phantomsection\label{GELFAND-0153-U801}\textbf{\texttt{GELFAND-0153-U801}} (citation-name).\nopagebreak[4]\par
В строке «В работе ... и Джеймса» первая фамилия имеет Т-образную начальную букву; далее читается «иппера».\par
Кандидаты: Типпера; Диппера. Чтение не установлено в дипломатическом исходнике.\par
Историческое примечание исходной транскрипции: Оригинал 2408×3448; прямоугольник (480,1690)-(860,1845), координаты от левого верхнего угла. Осмотрен оригинал и увеличения 3× и 4×. Начальная буква не заменена по библиографическому ожиданию; требуется сверка человеком.\par\medskip
\subsection*{\texttt{GELFAND-P001414}}
\par\noindent\leavevmode\phantomsection\label{GELFAND-0153-U901}\textbf{\texttt{GELFAND-0153-U901}} (eq11-cancelled-word).\nopagebreak[4]\par
Перед r\_n(t)=±1 в (11) зачёркнуто короткое слово; над ним вписано «Пусть» с указателем к месту замены. В зачёркнутом видны начальная петля и несколько соединённых штрихов.\par
Кандидаты: если. Чтение не установлено в дипломатическом исходнике.\par
Историческое примечание исходной транскрипции: Оригинал 2408×3448; (180,2650)-(360,2780). Просмотр 3× и 9×. «Пусть» уверенно читается как вставка; прежнее слово дано только как кандидат.\par\medskip
\par\noindent\leavevmode\phantomsection\label{GELFAND-0153-U902}\textbf{\texttt{GELFAND-0153-U902}} (eq13-terminal-mark).\nopagebreak[4]\par
После s,t∈T(λ), ниже строки (13), имеется отдельный небольшой двухштриховой знак с нижним хвостом.\par
Кандидаты: q; исправленная запятая. Чтение не установлено в дипломатическом исходнике.\par
Историческое примечание исходной транскрипции: Оригинал 2408×3448; (1020,3030)-(1130,3150). Просмотр 3×, 7× и 10×. Не установлено, является ли знак частью индекса, пунктуацией или исправлением; в формулу как достоверный индекс не включён.\par\medskip
\par\noindent\leavevmode\phantomsection\label{GELFAND-0153-U903}\textbf{\texttt{GELFAND-0153-U903}} (eq14-blotted-mark).\nopagebreak[4]\par
В (14) непосредственно после скобки аргумента первого множителя W виден плотный чёрный штрих ниже основной строки. Более ранняя форма закрыта чернилами.\par
Кандидаты: не записаны. Чтение не установлено в дипломатическом исходнике.\par
Историческое примечание исходной транскрипции: Оригинал 2408×3448; (785,3235)-(900,3320). Осмотр 5× и 8×. Неясно, зачёркнут ли математический знак или пунктуация; оба множителя и степень 1/2 читаются. След не удалён из дипломатического описания.\par\medskip
\par\noindent\leavevmode\phantomsection\label{GELFAND-0153-U904}\textbf{\texttt{GELFAND-0153-U904}} (eq14-cancelled-condition).\nopagebreak[4]\par
В начале последней строки зачёркнута фраза «при s ... t»; под отношением виден k. Отношение похоже на две волнистые линии, но перекрыто отменяющими штрихами.\par
Кандидаты: при s ≈\_k t; при s с другим знаком отношения с индексом k и t. Чтение не установлено в дипломатическом исходнике.\par
Историческое примечание исходной транскрипции: Оригинал 2408×3448; (245,3280)-(600,3420). Осмотр 4× и 6×. Вероятно, повтор условия из предыдущей строки; это лишь сопоставление, а не исправление источника. Зачёркнутый текст сохранён с вопросом.\par\medskip
\subsection*{Отдельное математическое замечание}\label{GELFAND-P001412-EDITORIAL-MATH-01}
Если сумму в (6) понимать как сумму по i=1,...,d (в рукописи пределы не указаны), то при d=2 и x1=y1=(1,0), x2=y2=(0,1) левая часть равна 1, а слагаемые справа равны -1 и 1. Это условная проверка несогласованности, а не исправление рукописи.\par
